% page 446 of the facsimile (image p-445 of the scan)
% block 149.9 x 242.0 mm ; left margin 8.0 mm
\pgc{-12.700mm}{0.000mm}{
\l{\cel{3}438.}
\l{}
\l{\sou{pikar}. (\sou{trans.}) I. Entamar neprofunde, per pinto. - II.}
\l{\cel{3}Truizar per ulo quan onu igas penetrar per pinto. -}
\l{\cel{3}III. Igar (ulo penetrar aden ulo, pinto-avane. - IV.}
\l{\cel{3}Stimular vivace ; frapar per sento vivaca, penetriva.}
\l{\cel{3}- DEFS.}
\l{}
\l{\sou{pikadoro}. Kavaliero qua, en tauro-kombato, atakas la ani-}
\l{\cel{3}malo per piquo. - DEFIRS.}
\l{}
\l{\sou{pikeo.}\cel{2}Stofo de kotono, ek du texuri aplikita ica kontre}
\l{\cel{3}ita, ed unionita per steburi qui formacas desegnuri.}
\l{\cel{3}- DEFIS.}
\l{}
\l{\sou{piketo}. I. (\sou{armeo)}.\cel{2}Mikra trupo de kavalriani, pronta de-}
\l{\cel{3}partar quik lor signalo. - II. Karto-ludo ek 32 karti,}
\l{\cel{3}en qua la ludanto qua ganis 30 punti ante ke la adverso}
\l{\cel{3}ganis un, kontas sua ganajo quale se lu atingabus 60.}
\l{\cel{3}DEFIRS.}
\l{}
\l{\sou{piklo}.- Singla ek la kondimenti vejetantala, konservita per}
\l{\cel{3}vinagro. - E.}
\l{}
\l{\sou{piknika}r. (\sou{netrans.}) Agar repasto en qua singla partopren-}
\l{\cel{3}anto pagas sua quoto di la spenso,od adportas sua dishi.}
\l{\cel{3}DEFRS.}
\l{}
\l{\sou{pikono}.\cel{2}Utensilo, implemento, ek fero, pintoza, quan uzas}
\l{\cel{3}la ministi karbonala, petrala, la terlaboristi, por ex-}
\l{\cel{3}kavar roko o metarii harda ; la fera parto hagas pinto}
\l{\cel{3}perpendikla ye la mancho. -}
\l{}
\l{\sou{piktar}. (\sou{trans}.) Reprezentar per farbi. - EFIS.}
\l{}
\l{\sou{pikverdo}.\cel{1}\sou{(zool.}) Ucelo kun plumaro flava e verda, ek la}
\l{\cel{3}genero \textquotedbl{}pegi\textquotedbl{}. - L.\cel{1}\sou{picus viridis}. - FIS.}
\l{}
\l{\sou{pilo.}\cel{2}(\sou{anat.)}\cel{3}Produkturo epidermala,qua havas la formo}
\l{\cel{3}di filo tenua, qua kovras la korpo di preske omna mami-}
\l{\cel{3}feri. - II. Ta produkturo che la homo, ube olu kreskas}
\l{\cel{3}sur ula parti di la korpo (vangi, labii, mentono,axeli,}
\l{\cel{3}pubio). - EFIS.}
\l{}
\l{\sou{pilastro}. (\sou{arkitekt.}) I. Kolono quadrata, maxim-multa-kaze}
\l{\cel{3}sinkita en muro. - II. Montanto quan onu fixigas inter-}
\l{\cel{3}valope en greto, pasamano, balkono. - III. Kolono de}
\l{\cel{3}petro-bloki, qua subtenas parti diversa en edifico, en}
\l{\cel{3}ponto, e c. - DEFIRS.}
\l{}
\l{\sou{pilgrimar}. (\sou{netrans.})\cel{2}Voyajar pro devoceso, pieso, a loko}
\l{\cel{3}sakrigita. - DEFIR.}
\l{}
\l{pilio. (\sou{elektro)}.\cel{2}Aparato (inventita da Volta, en la yaro}
\l{\cel{3}1800), qua transformas ad elektro-fluo la energio quan}
\l{\cel{3}developas reakto kemiala. - DEF.}
}
